\subsection{Transport de la section électronique vers le secteur électrofaible}
\label{amp-ew-seccion}

La section centrale de \eqref{md:hib:electron} conserve conjointement le
préfacteur, le correcteur et la mémoire de retour. Ses dépendances proviennent
des feuillets APP, de l’orientation TRIT et de la route enregistrée par TPK ;
le caractère angulaire et les tours agissent ensuite sur les multisections
de la même structure discrète du continu. Nous écrivons
\begin{equation}
 r_e=R_{\rm act}^{\kappa_e},\qquad
 E_e^\beta=\mathcal U_E\varepsilon_e^{(0)}r_e,\qquad
 \kappa_e=80-54\alpha+6\Delta_4.
 \label{amp-ew-electron}
\end{equation}
Ici $\mathcal U_E$ est la section de la droite énergétique et $r_e$ abrège
le retour électronique. La décade d’action reste dans sa droite propre ;
la coordonnée énergétique réalisée est transportée une seule fois.

Les signatures angulaires de \eqref{amp-higgs-firmas} produisent, dans la même
coordinatisation, $E_s^\angle=E_e^\beta\exp(n_sx+s_sy/6)$. En particulier,
\begin{equation}
 E_W^\angle=\mathcal U_E\varepsilon_e^{(0)}r_e
                  e^{94x-y/6},\qquad D=e^{-2x}.
 \label{amp-ew-w}
\end{equation}
Cette réalisation angulaire de W est utilisée. Les raffinements d’une autre tour
conservent leurs facteurs propres lorsqu’on compose avec eux.

\begin{proposition}[Échelle électronique et couplage de Fermi au niveau arbre]
Dans la réalisation électrofaible au niveau arbre, avec les coordonnées énergétiques
exprimées dans une même unité et $g^2=4\pi\alpha/(1-D)$, les relations
$v=2E_W^\angle/g$ et $G_F^{(0)}=1/(\sqrt2v^2)$ donnent
\begin{equation}
 G_F^{(0)}=
 \frac{\pi\alpha\,e^{-188x+y/3}}
 {\sqrt2(1-D)\mathcal U_E^2(\varepsilon_e^{(0)})^2r_e^2}.
 \label{amp-ew-fermi}
\end{equation}
Les signatures, les canaux et la coordinatisation énergétique étant fixes, le passage de
l’étape basale à l’étape après retour satisfait
\begin{equation}
 \frac{(G_F^{(0)})_\beta}{(G_F^{(0)})_0}=r_e^{-2},\qquad
 \frac{v_\beta}{v_0}=r_e.
 \label{amp-ew-escala}
\end{equation}
Les quotients $E_W^\angle/E_Z^\angle$, $E_H^\angle/E_W^\angle$ et
$\lambda_H=(g^2/8)(E_H^\angle/E_W^\angle)^2$ sont indépendants de
ce facteur commun $r_e$.
\end{proposition}
\begin{proof}
La substitution de $v$ donne $G_F^{(0)}=g^2/[4\sqrt2(E_W^\angle)^2]$.
Appliquer l’expression de $g^2$ et élever \eqref{amp-ew-w} au carré
produit \eqref{amp-ew-fermi} : les coefficients $188$ et $1/3$ résultent
respectivement de $2\cdot94$ et $2/6$. Remplacer $E_e^0$ par
$E_e^\beta=r_eE_e^0$ multiplie toutes les énergies angulaires par $r_e$.
La définition de $v$ conserve la première puissance ; celle de $G_F^{(0)}$
conserve l’inverse de son carré. Dans chaque quotient d’énergies, le
même facteur positif se simplifie, et la formule de $\lambda_H$ conserve
cette simplification. Les couplages $g,g'$ dépendent de $\alpha,x$ et
restent fixes dans cette comparaison.
\end{proof}

La grandeur $G_F^{(0)}$ a la dimension de l’inverse du carré d’une énergie.
La réalisation au niveau arbre déclarée spécifie la portée physique de cette
composition ; les corrections radiatives ont leurs opérateurs et leur ordre
postérieurs. Comparer les deux étapes du lecteur conserve leurs coordinatisations et
leurs signatures, au lieu de postuler une variation temporelle indépendante des
constantes corrélées.

\subsection{Conservation des raffinements de route}
\label{amp-ew-refinamiento}
La composition avec la loi complète préserve le caractère et la tour qui
correspondent à chaque route. Pour deux sections propres positives comparables,
dans une même coordinatisation et avec un même lecteur $r\in\{D,\Omega\}$, soit
\begin{equation}
 \rho_{W/e}^{\,r}
 =\frac{\chi_{\rm full}^{\,r}(\sigma_W,\eta_W)}
        {\chi_{\rm full}^{\,r}(\sigma_e,\eta_e)}
 R_{\rm act}^{\kappa_r(\Gamma_W)-\kappa_e}.
 \label{amp-ew-razon-completa}
\end{equation}
L’énergie correspondante est $E_W=E_e^\beta\rho_{W/e}^{\,r}$.
Dans la réalisation de Fermi compatible avec cette section, on obtient
\begin{equation}
 G_F^{(0)}=
 \frac{\pi\alpha}
 {\sqrt2(1-D)(E_e^\beta)^2(\rho_{W/e}^{\,r})^2}.
 \label{amp-ew-fermi-completo}
\end{equation}
La preuve substitue le rapport complet dans l’identité quadratique
précédente. Dans la coordinatisation angulaire explicite, ce rapport prend la forme de
\eqref{amp-ew-w} ; dans une réalisation raffinée, il retient également les
différences de tours et de lecteurs de \eqref{md:hib:cocientes}.
Si $E_s=E_e^\beta\chi_sf_s$, alors
$E_s/E_t=(\chi_s/\chi_t)(f_s/f_t)$. Cette égalité montre exactement
quel facteur se simplifie et quelle information spécifique demeure.

\subsection{Congruence massique et inversion spectrale sans commutation}
\label{amp-ew-congruencia}
Soit $\widehat M^0>0$ l’opérateur basal sur une fibre finie positive,
avec son caractère et ses tours incorporés, et soit $\widehat B^r$
le lecteur autoadjoint de retour. Le transport d’action produit
$C=R_{\rm act}^{\widehat B^r/2}>0$ et la loi
$\widehat M^\beta=C\widehat M^0C$.

\begin{proposition}[Réciprocité opératoire de la longueur d’action]
Avec $\hbar_a,c>0$ dans une même coordinatisation dimensionnelle,
\begin{equation}
 \begin{aligned}
 (\widehat M^\beta)^{-1}
    &=C^{-1}(\widehat M^0)^{-1}C^{-1},\\
 \widehat L_C^\beta:=\frac{\hbar_a}{c}(\widehat M^\beta)^{-1}
    &=C^{-1}\widehat L_C^0C^{-1},\qquad
 \widehat L_C^0=\frac{\hbar_a}{c}(\widehat M^0)^{-1}.
 \end{aligned}
 \label{amp-ew-inversa-operador}
\end{equation}
Ces identités admettent $[C,\widehat M^0]\ne0$.
\end{proposition}
\begin{proof}
La positivité de $C$ implique l’inversibilité. Pour $u\ne0$,
$\langle u,C\widehat M^0Cu\rangle
=\langle Cu,\widehat M^0Cu\rangle>0$ ; par conséquent, l’opérateur final
est également inversible. La multiplication directe dans les deux ordres donne
\[
 (C\widehat M^0C)\bigl(C^{-1}(\widehat M^0)^{-1}C^{-1}\bigr)=I,
 \qquad
 \bigl(C^{-1}(\widehat M^0)^{-1}C^{-1}\bigr)(C\widehat M^0C)=I.
\]
Multiplier par $\hbar_a/c$ prouve la seconde identité. Dans la
décomposition spectrale de l’opérateur final,
$\widehat M^\beta=\sum_jm_jP_j$ avec $m_j>0$, on obtient
$\widehat L_C^\beta=\sum_j\hbar_a/(cm_j)P_j$. La fonction inverse
est injective sur les valeurs propres positives : elle conserve leurs projecteurs
et leurs multiplicités finaux.
\end{proof}

La congruence peut changer les projecteurs de l’opérateur basal ;
l’inversion spectrale conserve ceux de l’opérateur final. Le secteur nul
reste en dehors de cette inversion. Une préparation qui occupe plusieurs
espaces propres conserve sa mesure spectrale et ses multiplicités.
En rang un, l’identité retrouve la longueur électronique
$\bar\lambda_{C,e}=\hbar_ac/E_e^\beta$ et, ensuite,
$a_{{\rm B},e}=\bar\lambda_{C,e}/\alpha$. Le résultat relie
l’échelle électronique, la réponse électrofaible et la longueur d’action
par des opérations explicites sur leurs structures partagées.
